\documentclass[11pt]{article}
\usepackage[margin=1in]{geometry}
\usepackage{hyperref}
\usepackage{enumitem}
\usepackage{booktabs}
\usepackage{xcolor}
\title{ProbOS --- Month 2, Week 7\\FastAPI Service Layer}
\author{Nisong Monyimba}
\date{\today}

\begin{document}
\maketitle

\section*{Week 7 goal}
Expose the kernel over HTTP so it is usable outside a Python REPL --
the first step toward external/enterprise-facing usage referenced in
\texttt{docs/monthly\_plans/overall/main.tex}.

\section*{Day-by-day plan}

\subsection*{Day 1 -- FastAPI scaffold}
\begin{itemize}[leftmargin=*]
  \item Set up \texttt{python/server/} package structure
  \item Basic FastAPI app, health check endpoint
  \item Pydantic schemas mirroring \texttt{Model}/\texttt{Distribution}
    constructor signatures
\end{itemize}

\subsection*{Day 2 -- /simulate endpoint}
\begin{itemize}[leftmargin=*]
  \item \texttt{POST /simulate}: accept model name, priors, $N$,
    \texttt{n\_steps}; run \texttt{MonteCarloEngine}; return
    percentiles + convergence summary as JSON
  \item Input validation: sane upper bounds on $N$/\texttt{n\_steps} to
    prevent resource exhaustion (per
    \texttt{docs/standards/quality\_standards.md} Section 8, deferred
    item now becoming relevant)
\end{itemize}

\subsection*{Day 3 -- /sensitivity endpoint}
\begin{itemize}[leftmargin=*]
  \item \texttt{POST /sensitivity}: run \texttt{SobolSensitivity},
    return $S_1$/$S_T$ indices as JSON
\end{itemize}

\subsection*{Day 4 -- /filter endpoint}
\begin{itemize}[leftmargin=*]
  \item \texttt{POST /filter}: run \texttt{ParticleFilter} against
    posted observation data, return posterior mean/std trajectory
\end{itemize}

\subsection*{Day 5 -- Integration tests}
\begin{itemize}[leftmargin=*]
  \item \texttt{python/tests/test\_server.py}: FastAPI
    \texttt{TestClient}-based tests for all endpoints
  \item Confirm JSON responses match direct Python kernel calls
    numerically
\end{itemize}

\subsection*{Day 6 -- Pre-week audit + CI integration}
\begin{itemize}[leftmargin=*]
  \item Run \texttt{pre\_week\_audit.sh} before considering the week
    complete
  \item Extend CI to install and test the server layer
  \item Extend \texttt{docs/standards/quality\_standards.md} Section 8
    with the now-relevant API/service checks (previously listed as
    deferred)
\end{itemize}

\subsection*{Day 7 -- Study guide + documentation}
\begin{itemize}[leftmargin=*]
  \item \texttt{docs/study/study\_guide.md} entries: FastAPI
    documentation, Pydantic validation patterns
  \item Week 7 retrospective appended to this document
\end{itemize}

\section*{Exit criteria}
\texttt{curl -X POST localhost:8000/simulate ...} returns valid JSON
matching a direct Python call to the same \texttt{MonteCarloEngine},
for all three endpoints (\texttt{/simulate}, \texttt{/sensitivity},
\texttt{/filter}).


\section*{What was actually accomplished (retrospective)}

\begin{itemize}[leftmargin=*]
  \item \texttt{python/server/schemas.py}: Pydantic request/response
    schemas for all four endpoints, with resource-exhaustion bounds
    enforced via \texttt{Field(ge=..., le=...)} at the HTTP boundary
    -- exactly the previously-deferred check from
    \texttt{quality\_standards.md} Section 8, now active
  \item \texttt{python/server/main.py}: \texttt{GET /health},
    \texttt{POST /simulate}, \texttt{POST /sensitivity},
    \texttt{POST /filter} -- each a thin translation layer over an
    existing kernel object, no new kernel logic
  \item Kernel-raised \texttt{ValueError} surfaces as HTTP 422 (not an
    unhandled 500); unknown \texttt{model\_name} surfaces as HTTP 404
  \item 19 integration tests, including three that directly compare
    HTTP JSON responses against a matching direct Python kernel call
    at \texttt{rtol=1e-10} -- proving zero numerical discrepancy
    introduced by the service layer, not just ``returns 200''
  \item Additionally smoke-tested with a real \texttt{uvicorn} process
    over an actual HTTP socket (\texttt{/health} and \texttt{/simulate}),
    not just FastAPI's in-process \texttt{TestClient}
  \item \texttt{quality\_standards.md} Section 8 updated: the
    previously-deferred API/service checks are now active and checked
    off, each pointing to the specific test that verifies it
  \item CI extended across all three relevant jobs
    (\texttt{python-tests}, \texttt{integration}, \texttt{quality-audit})
    to install and test the new server layer
\end{itemize}

\section*{Numbers at Week 7 close}
\begin{itemize}[leftmargin=*]
  \item Python tests: 341/341 passing (322 + 19 new)
  \item C++ tests: 13/13 passing
  \item mypy strict: 0 errors (full \texttt{python/} tree including
    \texttt{python/server/})
  \item ruff: 0 warnings
  \item \texttt{pre\_week\_audit.sh}: 21/21 checks passed
\end{itemize}

\section*{Carried into later weeks}
Only the \texttt{battery} model is currently registered in the server
layer's model registry. Extending it to the Week 4 cross-discipline
models (option pricer, ED queue, clinical trial) is a natural, low-risk
follow-up once genuinely needed. The service layer itself is complete
and satisfies its exit criteria exactly as planned.
\end{document}
